\documentclass[../../main.tex]{subfiles}

\graphicspath{{../../image/}{../../figure/}}

\begin{document}

\section{习题}\label{section:ch03-习题}

\begin{exercise}\label{exercise:ch03-01}
（★）验证 Bernoulli 分布[\eqref{align:ch03-Bernoulli分布}]满足下列性质：
\begin{align*}
\sum_{x=0}^{1}p(x\mid\mu)=1,\qquad \mathbb{E}[x]=\mu,\qquad \mathrm{var}[x]=\mu(1-\mu),
\end{align*}
并证明符合 Bernoulli 分布的随机二元变量 $x$ 的熵 $H[x]$ 由下式给出：
\begin{align*}
H[x]=-\mu\ln\mu-(1-\mu)\ln(1-\mu).
\end{align*}
\end{exercise}

\begin{exercise}\label{exercise:ch03-02}
（★★）\eqref{align:ch03-Bernoulli分布}给出的 Bernoulli 分布形式在两个 $x$ 值之间不对称。在某些情况下，使用另一种等价的形式会更方便，其中 $x\in\{-1,1\}$，此时该分布可以写成
\begin{align*}
p(x\mid\mu)=\left(\frac{1-\mu}{2}\right)^{(1-x)/2}\left(\frac{1+\mu}{2}\right)^{(1+x)/2},
\end{align*}
其中 $\mu\in[-1,1]$。证明该式给出的分布是归一化的，并计算它的均值、方差和熵。
\end{exercise}

\begin{exercise}\label{exercise:ch03-03}
（★★）证明二项分布[\eqref{align:ch03-二项分布}]是归一化的。为此，首先使用定义\eqref{align:ch03-组合数}，即从总数为 $N$ 的物体中选择 $m$ 个相同物体的组合数量，证明下式：
\begin{align*}
\binom{N}{m}+\binom{N}{m-1}=\binom{N+1}{m}.
\end{align*}
然后利用上述结论，通过归纳证明以下结论：
\begin{align*}
(1+x)^{N}=\sum_{m=0}^{N}\binom{N}{m}x^{m},
\end{align*}
这个称为二项式定理（binomial theorem）的结论对所有实数 $x$ 都成立。最后，证明二项分布是归一化的，从而使
\begin{align*}
\sum_{m=0}^{N}\binom{N}{m}\mu^{m}(1-\mu)^{N-m}=1.
\end{align*}
\end{exercise}

\begin{exercise}\label{exercise:ch03-04}
（★★）证明二项分布的均值由\eqref{align:ch03-二项均值}给出。为此，在归一化条件
\begin{align*}
\sum_{m=0}^{N}\binom{N}{m}\mu^{m}(1-\mu)^{N-m}=1
\end{align*}
的两边对 $\mu$ 求导，并重排以获得 $\mu$ 的均值表达式。类似地，通过对上式关于 $\mu$ 进行二次求导，并利用二项分布均值的结果，证明二项分布方差的结果[\eqref{align:ch03-二项方差}]。
\end{exercise}

\begin{exercise}\label{exercise:ch03-05}
（★）证明多元 Gauss 分布\eqref{align:ch03-多元Gauss}的众数为 $\boldsymbol{\mu}$。
\end{exercise}

\begin{exercise}\label{exercise:ch03-06}
（★★）假设 $\boldsymbol{x}$ 服从均值为 $\boldsymbol{\mu}$、协方差为 $\boldsymbol{\Sigma}$ 的 Gauss 分布。证明线性变换后的变量 $\boldsymbol{A}\boldsymbol{x}+\boldsymbol{b}$ 也服从 Gauss 分布，并给出它的均值和协方差。
\end{exercise}

\begin{exercise}\label{exercise:ch03-07}
（★★★）证明两个 Gauss 分布 $q(\boldsymbol{x})=\mathcal{N}(\boldsymbol{x}\mid\boldsymbol{\mu}_{q},\boldsymbol{\Sigma}_{q})$ 和 $p(\boldsymbol{x})=\mathcal{N}(\boldsymbol{x}\mid\boldsymbol{\mu}_{p},\boldsymbol{\Sigma}_{p})$ 之间的 Kullback-Leibler 散度为
\begin{align*}
\mathrm{KL}(q\,\|\,p)=\frac{1}{2}\left\{\ln\frac{\lvert\boldsymbol{\Sigma}_{p}\rvert}{\lvert\boldsymbol{\Sigma}_{q}\rvert}-D+\mathrm{tr}\bigl(\boldsymbol{\Sigma}_{p}^{-1}\boldsymbol{\Sigma}_{q}\bigr)+(\boldsymbol{\mu}_{p}-\boldsymbol{\mu}_{q})^{\mathsf{T}}\boldsymbol{\Sigma}_{p}^{-1}(\boldsymbol{\mu}_{p}-\boldsymbol{\mu}_{q})\right\},
\end{align*}
其中 $\mathrm{tr}(\cdot)$ 表示矩阵的迹，$D$ 为 $\boldsymbol{x}$ 的维度。
\end{exercise}

\begin{exercise}\label{exercise:ch03-08}
（★★）这个习题表明，对于给定的协方差，具有最大熵的多元分布是 Gauss 分布。分布 $p(\boldsymbol{x})$ 的熵由下式给出：
\begin{align*}
H[\boldsymbol{x}]=-\int p(\boldsymbol{x})\ln p(\boldsymbol{x})\,\mathrm{d}\boldsymbol{x}.
\end{align*}
我们希望在所有分布 $p(\boldsymbol{x})$ 中最大化 $H[\boldsymbol{x}]$，约束如下：满足 $p(\boldsymbol{x})$ 是归一化的，且具有特定的均值和协方差，使得
\begin{align*}
\int p(\boldsymbol{x})\,\mathrm{d}\boldsymbol{x}=1,\qquad \int p(\boldsymbol{x})\boldsymbol{x}\,\mathrm{d}\boldsymbol{x}=\boldsymbol{\mu},\qquad \int p(\boldsymbol{x})(\boldsymbol{x}-\boldsymbol{\mu})(\boldsymbol{x}-\boldsymbol{\mu})^{\mathsf{T}}\mathrm{d}\boldsymbol{x}=\boldsymbol{\Sigma}.
\end{align*}
通过对熵进行变分最大化，并使用 Lagrange 乘子来确保约束条件，证明最大熵分布由 Gauss 分布\eqref{align:ch03-多元Gauss}给出。
\end{exercise}

\begin{exercise}\label{exercise:ch03-09}
（★★★）证明多元 Gauss 分布 $\mathcal{N}(\boldsymbol{x}\mid\boldsymbol{\mu},\boldsymbol{\Sigma})$ 的熵由下式给出：
\begin{align*}
H[\boldsymbol{x}]=\frac{1}{2}\ln\lvert\boldsymbol{\Sigma}\rvert+\frac{D}{2}\bigl(1+\ln(2\pi)\bigr),
\end{align*}
其中 $D$ 为 $\boldsymbol{x}$ 的维度。
\end{exercise}

\begin{exercise}\label{exercise:ch03-10}
（★★★）考虑两个随机变量 $x_1$ 和 $x_2$，它们的分布都是 Gauss 分布，均值分别为 $\mu_1$ 和 $\mu_2$，精度分别为 $\tau_1$ 和 $\tau_2$。推导变量 $x=x_1+x_2$ 的微分熵表达式。为此，首先利用以下关系找到 $x$ 的分布：
\begin{align*}
p(x)=\int_{-\infty}^{\infty}p(x\mid x_2)p(x_2)\,\mathrm{d}x_2,
\end{align*}
并完成指数中的平方项。然后观察到这表示两个 Gauss 分布的卷积，而卷积本身也是 Gauss 分布。最后利用式\eqref{align:ch03-一元Gauss}中一元 Gauss 分布的熵的计算结果。
\end{exercise}

\begin{exercise}\label{exercise:ch03-11}
（★）考虑\eqref{align:ch03-多元Gauss}给出的多元 Gauss 分布。通过将精度矩阵（逆协方差矩阵）写成对称矩阵和反对称矩阵的和，证明反对称项不会出现在 Gauss 分布的指数中，所以精度矩阵可以不失一般性地假定为对称的。又因为对称矩阵的逆矩阵也是对称的（见\cref{exercise:ch03-16}），所以协方差矩阵也可以不失一般性地选为对称的。
\end{exercise}

\begin{exercise}\label{exercise:ch03-12}
（★★★）考虑一个实对称矩阵 $\boldsymbol{\Sigma}$，其特征方程由\eqref{align:ch03-特征方程}给出。通过取这个方程的复共轭，并减去原方程，然后计算与特征向量 $\boldsymbol{u}_i$ 的内积，证明特征值 $\lambda_i$ 是实数。类似地，利用 $\boldsymbol{\Sigma}$ 的对称性，证明当 $\lambda_i\neq\lambda_j$ 时，特征向量 $\boldsymbol{u}_i$ 和 $\boldsymbol{u}_j$ 是正交的。最后，请证明特征向量集可以不失一般性地选择为正交的，使得它们满足\eqref{align:ch03-特征向量正交}，即使一些特征值为零。
\end{exercise}

\begin{exercise}\label{exercise:ch03-13}
（★）证明\eqref{align:ch03-特征方程}这样一个具有特征向量方程的实对称矩阵 $\boldsymbol{\Sigma}$ 可以表示为形如\eqref{align:ch03-协方差谱分解}的特征向量的展开形式，且系数由特征值给出。类似地，证明逆矩阵 $\boldsymbol{\Sigma}^{-1}$ 具有形如\eqref{align:ch03-逆协方差谱分解}的表示形式。
\end{exercise}

\begin{exercise}\label{exercise:ch03-14}
（★★）正定矩阵 $\boldsymbol{\Sigma}$ 可以定义如下：对于任意实向量 $\boldsymbol{a}$，其二次型
\begin{align*}
\boldsymbol{a}^{\mathsf{T}}\boldsymbol{\Sigma}\boldsymbol{a}
\end{align*}
是正的。证明 $\boldsymbol{\Sigma}$ 为正定的一个充要条件是，$\boldsymbol{\Sigma}$ 的所有特征值 $\lambda_i$ 都是正的，其中特征值由\eqref{align:ch03-特征方程}定义。
\end{exercise}

\begin{exercise}\label{exercise:ch03-15}
（★）证明一个大小为 $D\times D$ 的实对称矩阵具有 $D(D+1)/2$ 个独立的参数。
\end{exercise}

\begin{exercise}\label{exercise:ch03-16}
（★）证明对称矩阵的逆矩阵也是对称的。
\end{exercise}

\begin{exercise}\label{exercise:ch03-17}
（★★）通过使用特征向量展开式\eqref{align:ch03-协方差谱分解}对坐标系进行对角化，证明与常数 Mahalanobis 距离 $\Delta$ 对应的超椭球的体积为
\begin{align*}
V_{D}\,\lvert\boldsymbol{\Sigma}\rvert^{1/2}\,\Delta^{D},
\end{align*}
其中 $V_{D}$ 是 $D$ 维单位球的体积，而 Mahalanobis 距离由\eqref{align:ch03-马氏距离}定义。
\end{exercise}

\begin{exercise}\label{exercise:ch03-18}
（★★）利用定义\eqref{align:ch03-舒尔补}，并通过对它的两边乘以下列矩阵，证明恒等式\eqref{align:ch03-分块逆}成立。
\begin{align*}
\begin{pmatrix}\boldsymbol{A} & \boldsymbol{B}\\ \boldsymbol{C} & \boldsymbol{D}\end{pmatrix}
\end{align*}
\end{exercise}

\begin{exercise}\label{exercise:ch03-19}
（★★★）\cref{subsection:ch03-条件分布}和\cref{subsection:ch03-边缘分布}介绍了多元 Gauss 分布的条件分布和边缘分布。更一般地，可以将 $\boldsymbol{x}$ 的分量划分为三组——$\boldsymbol{x}_{a}$、$\boldsymbol{x}_{b}$ 和 $\boldsymbol{x}_{c}$，对应地，将均值向量 $\boldsymbol{\mu}$ 和协方差矩阵 $\boldsymbol{\Sigma}$ 也按照如下形式进行划分：
\begin{align*}
\boldsymbol{\mu}=\begin{pmatrix}\boldsymbol{\mu}_{a}\\ \boldsymbol{\mu}_{b}\\ \boldsymbol{\mu}_{c}\end{pmatrix},\qquad
\boldsymbol{\Sigma}=\begin{pmatrix}\boldsymbol{\Sigma}_{aa} & \boldsymbol{\Sigma}_{ab} & \boldsymbol{\Sigma}_{ac}\\ \boldsymbol{\Sigma}_{ba} & \boldsymbol{\Sigma}_{bb} & \boldsymbol{\Sigma}_{bc}\\ \boldsymbol{\Sigma}_{ca} & \boldsymbol{\Sigma}_{cb} & \boldsymbol{\Sigma}_{cc}\end{pmatrix}.
\end{align*}
利用\cref{section:ch03-多元Gauss分布}的结论，给出条件分布 $p(\boldsymbol{x}_{a}\mid\boldsymbol{x}_{b})$ 的表达式，其中 $\boldsymbol{x}_{c}$ 已被边缘化。
\end{exercise}

\begin{exercise}\label{exercise:ch03-20}
（★★）线性代数中一个非常有用的结论是 Woodbury 矩阵求逆公式，即
\begin{align*}
(\boldsymbol{A}+\boldsymbol{B}\boldsymbol{C}\boldsymbol{D})^{-1}=\boldsymbol{A}^{-1}-\boldsymbol{A}^{-1}\boldsymbol{B}(\boldsymbol{C}^{-1}+\boldsymbol{D}\boldsymbol{A}^{-1}\boldsymbol{B})^{-1}\boldsymbol{D}\boldsymbol{A}^{-1}.
\end{align*}
通过对\eqref{align:ch03-分块逆}的两边乘以 $(\boldsymbol{A}+\boldsymbol{B}\boldsymbol{C}\boldsymbol{D})$，证明这个结论的正确性。
\end{exercise}

\begin{exercise}\label{exercise:ch03-21}
（★）假设 $\boldsymbol{x}$ 和 $\boldsymbol{z}$ 是两个独立的随机向量，故 $p(\boldsymbol{x},\boldsymbol{z})=p(\boldsymbol{x})p(\boldsymbol{z})$。证明这两个随机向量的和 $\boldsymbol{y}=\boldsymbol{x}+\boldsymbol{z}$ 的均值等于其中每个随机向量的均值之和。类似地，证明 $\boldsymbol{y}$ 的协方差矩阵等于 $\boldsymbol{x}$ 与 $\boldsymbol{z}$ 的协方差矩阵之和。
\end{exercise}

\begin{exercise}\label{exercise:ch03-22}
（★★★）考虑变量
\begin{align*}
\boldsymbol{z}=\begin{pmatrix}\boldsymbol{x}\\ \boldsymbol{y}\end{pmatrix}
\end{align*}
的联合分布，它的均值和协方差分别由\eqref{align:ch03-联合均值化简}和\eqref{align:ch03-联合协方差}给出。利用\eqref{align:ch03-边缘均值}和\eqref{align:ch03-边缘协方差}，证明边缘分布 $p(\boldsymbol{x})$ 由\eqref{align:ch03-先验Gauss}给出。类似地，利用\eqref{align:ch03-条件均值}和\eqref{align:ch03-条件协方差}，证明条件分布 $p(\boldsymbol{y}\mid\boldsymbol{x})$ 由\eqref{align:ch03-线性Gauss条件}给出。
\end{exercise}

\begin{exercise}\label{exercise:ch03-23}
（★★）请利用分块矩阵求逆公式[\eqref{align:ch03-分块逆}]，证明精度矩阵[\eqref{align:ch03-联合精度}]的逆矩阵由协方差矩阵[\eqref{align:ch03-联合协方差}]给出。
\end{exercise}

\begin{exercise}\label{exercise:ch03-24}
（★★）从\eqref{align:ch03-联合均值}开始，利用\eqref{align:ch03-联合协方差}，验证\eqref{align:ch03-联合均值化简}。
\end{exercise}

\begin{exercise}\label{exercise:ch03-25}
（★★）考虑两个多维随机向量 $\boldsymbol{x}$ 和 $\boldsymbol{z}$，它们分别服从 Gauss 分布 $p(\boldsymbol{x})=\mathcal{N}(\boldsymbol{x}\mid\boldsymbol{\mu}_{x},\boldsymbol{\Sigma}_{x})$ 和 $p(\boldsymbol{z})=\mathcal{N}(\boldsymbol{z}\mid\boldsymbol{\mu}_{z},\boldsymbol{\Sigma}_{z})$，定义 $\boldsymbol{y}=\boldsymbol{x}+\boldsymbol{z}$。通过考虑由边缘分布 $p(\boldsymbol{x})$ 和条件分布 $p(\boldsymbol{y}\mid\boldsymbol{x})$ 的乘积组成的线性-Gauss 模型，并利用\eqref{align:ch03-边缘Bayes}和\eqref{align:ch03-边缘Bayes协方差}，证明边缘分布 $p(\boldsymbol{y})$ 由下式给出：
\begin{align*}
p(\boldsymbol{y})=\mathcal{N}(\boldsymbol{y}\mid\boldsymbol{\mu}_{x}+\boldsymbol{\mu}_{z},\boldsymbol{\Sigma}_{x}+\boldsymbol{\Sigma}_{z}).
\end{align*}
\end{exercise}

\begin{exercise}\label{exercise:ch03-26}
（★★★）本习题和下一习题提供了关于操作线性-Gauss 模型中出现的二次型的实践，它们也可以作为对正文中结论的独立检查。考虑由\eqref{align:ch03-先验Gauss}和\eqref{align:ch03-线性Gauss条件}给出的边缘分布和条件分布定义的联合分布 $p(\boldsymbol{x},\boldsymbol{y})$。通过检查联合分布指数中的二次型，并使用\cref{subsection:ch03-条件分布}讨论的“完成平方”技术，给出边缘分布 $p(\boldsymbol{y})$ 的均值表达式和协方差表达式，其中变量 $\boldsymbol{x}$ 已经被积分消除。为此，利用 Woodbury 矩阵求逆公式[\eqref{align:ch03-分块逆}]，验证这些结果与\eqref{align:ch03-边缘Bayes}和\eqref{align:ch03-边缘Bayes协方差}的一致性。
\end{exercise}

\begin{exercise}\label{exercise:ch03-27}
（★★★）考虑与\cref{exercise:ch03-26}中相同的联合分布，但使用“完成平方”技术来找到条件分布 $p(\boldsymbol{x}\mid\boldsymbol{y})$ 的均值表达式和协方差表达式。再次验证这些结果与\eqref{align:ch03-后验均值}和\eqref{align:ch03-后验协方差}的一致性。
\end{exercise}

\begin{exercise}\label{exercise:ch03-28}
（★★）为了找到多元 Gauss 分布的协方差矩阵的最大似然解，需要最大化对于 $\boldsymbol{\Sigma}$ 的对数似然函数[\eqref{align:ch03-Gauss对数似然}]，注意协方差矩阵必须是对称且正定的。此处忽略这些约束，直接进行最大化。利用\cref{section:appA-矩阵导数}中的矩阵求导法则，证明使对数似然函数[\eqref{align:ch03-Gauss对数似然}]最大化的协方差矩阵 $\boldsymbol{\Sigma}$ 是由样本协方差[\eqref{align:ch03-Gauss协方差ML}]给出的。注意最终的结果必然是对称且正定的（假设样本协方差是非奇异的）。
\end{exercise}

\begin{exercise}\label{exercise:ch03-29}
（★★）使用结果\eqref{align:ch03-Gauss均值}证明\eqref{align:ch03-二阶矩}。然后，利用\eqref{align:ch03-Gauss均值}和\eqref{align:ch03-二阶矩}证明
\begin{align*}
\mathbb{E}\bigl[\boldsymbol{x}_n\boldsymbol{x}_m^{\mathsf{T}}\bigr]=\boldsymbol{\mu}\boldsymbol{\mu}^{\mathsf{T}}+I_{nm}\boldsymbol{\Sigma},
\end{align*}
其中 $\boldsymbol{x}_n$ 表示从均值为 $\boldsymbol{\mu}$、协方差为 $\boldsymbol{\Sigma}$ 的 Gauss 分布中采样得到的数据点，而 $I_{nm}$ 表示单位矩阵中位于 $(n,m)$ 的元素。最后证明\eqref{align:ch03-协方差ML期望}。
\end{exercise}

\begin{exercise}\label{exercise:ch03-30}
（★）本章在讨论周期变量时使用的各种三角恒等式可以很容易地从下面的关系中证明：
\begin{align*}
\exp(\mathrm{i}A)=\cos A+\mathrm{i}\sin A,
\end{align*}
其中 $\mathrm{i}=\sqrt{-1}$。通过考虑恒等式 $\exp(\mathrm{i}A)\exp(-\mathrm{i}A)=1$ 证明\eqref{align:ch03-三角恒等式1}。类似地，使用恒等式 $\cos(A-B)=\Re\exp\{\mathrm{i}(A-B)\}$ 证明\eqref{align:ch03-三角恒等式2}，其中 $\Re$ 表示实部。最后，使用 $\sin(A-B)=\Im\exp\{\mathrm{i}(A-B)\}$（其中 $\Im$ 表示虚部）证明\eqref{align:ch03-正弦差}。
\end{exercise}

\begin{exercise}\label{exercise:ch03-31}
（★★）对于较大的 $m$，von Mises 分布[\eqref{align:ch03-冯米塞斯分布}]将会在众数 $\theta_0$ 处变得十分尖锐。通过定义 $\xi=m^{1/2}(\theta-\theta_0)$，并利用余弦函数的 Taylor 展开式，即
\begin{align*}
\cos\alpha=1-\frac{\alpha^{2}}{2}+O(\alpha^{4}),
\end{align*}
证明当 $m\to\infty$ 时，von Mises 分布趋向于 Gauss 分布。
\end{exercise}

\begin{exercise}\label{exercise:ch03-32}
（★）利用三角恒等式\eqref{align:ch03-正弦差}，证明\eqref{align:ch03-冯米塞斯驻点}对于 $\theta_0$ 的解由\eqref{align:ch03-冯米塞斯均值ML}给出。
\end{exercise}

\begin{exercise}\label{exercise:ch03-33}
（★）通过计算\eqref{align:ch03-冯米塞斯分布}给出的 von Mises 分布的一阶导数和二阶导数，并利用 $I_0(m)>0$（$m>0$），证明该分布的极大值发生在 $\theta=\theta_0$ 处，极小值发生在 $\theta=\theta_0+\pi\ (\mathrm{mod}\,2\pi)$ 处。
\end{exercise}

\begin{exercise}\label{exercise:ch03-34}
（★）通过利用结果\eqref{align:ch03-周期分量}，并结合\eqref{align:ch03-冯米塞斯均值ML}和三角恒等式\eqref{align:ch03-三角恒等式2}，证明 von Mises 分布的聚焦参数的最大似然解 $m_{\mathrm{ML}}$ 满足 $A(m_{\mathrm{ML}})=\bar{r}$，其中 $\bar{r}$ 是观测值的均值的半径。如\cref{figure:ch03-圆周均值}所示，它可以视为二维欧氏平面上的单位向量。
\end{exercise}

\begin{exercise}\label{exercise:ch03-35}
（★）验证多元 Gauss 分布可以表示为指数族分布的形式[\eqref{align:ch03-指数族定义}]，并为 $\boldsymbol{\eta}$、$\boldsymbol{u}(\boldsymbol{x})$、$h(\boldsymbol{x})$ 和 $g(\boldsymbol{\eta})$ 推导类似于\eqref{align:ch03-Gauss自然参数}至\eqref{align:ch03-Gauss归一化}的表达式。
\end{exercise}

\begin{exercise}\label{exercise:ch03-36}
（★）\eqref{align:ch03-指数族矩}表明，指数族分布中 $\ln g(\boldsymbol{\eta})$ 的负梯度由 $\boldsymbol{u}(\boldsymbol{x})$ 的期望给出。通过对\eqref{align:ch03-指数族归一}求二阶导数，证明
\begin{align*}
-\nabla\nabla\ln g(\boldsymbol{\eta})=\mathbb{E}\bigl[\boldsymbol{u}(\boldsymbol{x})\boldsymbol{u}(\boldsymbol{x})^{\mathsf{T}}\bigr]-\mathbb{E}[\boldsymbol{u}(\boldsymbol{x})]\,\mathbb{E}[\boldsymbol{u}(\boldsymbol{x})]^{\mathsf{T}}=\mathrm{cov}[\boldsymbol{u}(\boldsymbol{x})].
\end{align*}
\end{exercise}

\begin{exercise}\label{exercise:ch03-37}
（★★）考虑一个类似直方图的密度模型，其中空间 $\boldsymbol{x}$ 被划分为固定的区域，密度 $p(\boldsymbol{x})$ 在第 $i$ 个区域上取常数值 $h_i$。第 $i$ 个区域的体积表示为 $\Delta_i$。假设有一个集合，里面包含 $N$ 个 $\boldsymbol{x}$ 的观测值，其中 $n_i$ 个观测值落在第 $i$ 个区域。利用 Lagrange 乘子确保密度的归一化约束，并推导 $\{h_i\}$ 的最大似然估计量的表达式。
\end{exercise}

\begin{exercise}\label{exercise:ch03-38}
（★）证明 $K$ 近邻密度模型定义了一个反常分布，其在整个空间上的积分是发散的。
\end{exercise}

\end{document}
